% E13 — Benefit-Aware controller / uplift-CATE. The paper's sharpest open-problem
% statement. ZERO-GPU on the per-seed E6 cache. Fig: results/figures/e13_benefit_aware.png
% Table: results/tables/e13_benefit_aware.json.
% HONESTY: CATE does NOT beat random (significantly negative on gsm8k); the finding is
% that the treatment effect is UNPREDICTABLE from round-0 features despite a high oracle
% ceiling. Do not spin this as a working controller.

\subsection{Failure prediction is not benefit prediction}
\label{sec:rq13}

The hazard-guided controller (\S\ref{sec:rq2main}; full results in App.~\ref{app:controller}) targets
the E3a \emph{hazard} --- the probability that a debate fails. But what a controller should target is
the \emph{benefit} of intervening, and its robustness battery (App.~\ref{app:controller}) already
hinted the two diverge: the hazard ranks failures, yet targeting it does not beat a random budget. Here we make the distinction
exact. Because the E6 precompute stored \emph{both} counterfactual arms for every debate,
we observe both potential outcomes and therefore the individual treatment effect
$\mathrm{ITE}_i = Y_1^{(i)}-Y_0^{(i)}\in\{-1,0,+1\}$ ($Y_1$: combined-arm recovery, $Y_0$:
control-arm recovery). Any policy intervening on a set $S$ has accuracy
$\mathrm{never}+\tfrac{1}{n}\sum_{i\in S}\mathrm{ITE}_i$, so the optimal policy targets the
largest ITE, and targeting the \emph{true} ITE is an oracle ceiling. We learn the CATE
$\hat\tau(x)=\mathbb{E}[\mathrm{ITE}\mid X{=}x]$ from round-0 features (both a direct ITE
regressor and a two-model T-learner, group-cross-fitted over questions) and compare
policies at matched budget (Figure~\ref{fig:e13}).

Two facts stand out. \emph{First, the ceiling is high}: an oracle targeting true ITE reaches
always-on accuracy at a $\sim\!20\%$ budget, a gain over never of $+16.5$~points pooled
($+11.8$ on GSM8K, $+20.4$ on BBH) --- there is a great deal of benefit to capture.
\emph{Second, none of it is reachable from round-0 features.} The control outcome $Y_0$ is
predictable (OOF AUC $0.73$ --- this is essentially the hazard) and $Y_1$ is predictable
(AUC $0.67$), but their \emph{difference} is not: predicting whether intervention has
\emph{any} effect is barely above chance (AUC $0.58$) and predicting its \emph{sign} fails
outright --- the learned CATE ranks anti-correlated with the true ITE (Spearman $\le 0$ for
both learners, all scopes). Consequently the Benefit-Aware controller does \emph{not} beat a
random budget: pooled $\mathrm{AUC}_{20\text{-}40}$ of CATE$-$random is $-0.003$
(CI $[-0.023,+0.017]$), and on GSM8K it is significantly \emph{negative} ($-0.033$,
CI $[-0.060,-0.006]$), where the anti-correlation actively misleads. It captures
essentially $0\%$ of the oracle gain.

This is the paper's sharpest open-problem statement: \emph{failure prediction and
intervention-benefit prediction are different problems, and only the first is solved from
cheap transcript features.} Differencing two individually-predictable outcomes leaves a
treatment effect that the round-0 behavioral state does not determine --- which is exactly
why the hazard-guided controller matches, but cannot beat, a random budget (App.~\ref{app:controller}).
The large oracle ceiling shows the payoff is real; realizing it needs a benefit signal the
pre-intervention transcript does not carry --- semantic content, mid-debate dynamics, or a
cheap probe of the intervention itself. We therefore position the Benefit-Aware controller
as the correct formulation and its feature side as the central open problem, rather than
overselling a controller that only matches always-on.

\paragraph{Two pre-registered rescue attempts fail (E14, E15).}
One could object that the failure is an artifact of cheap handcrafted features, or of
using only \emph{passive} information. We ran two pre-registered rescues, each with a gate
fixed in advance (predictive: $\mathrm{AUC}_{\text{effect}}>0.65$ or
$\mathrm{Spearman}(\hat\tau,\mathrm{ITE})>0.20$; decision: a cluster-bootstrap CI on
benefit$-$random excluding zero in the $20$--$40\%$ band), same ITE, same learners, no
model search. \textbf{E14} changes \emph{only} the feature vector to residualized
\textsc{nomic-embed} semantic features of each agent's round-0 reasoning (own content, peer
similarity, own-vs-peer gap; per-question centroid removed to kill topic leakage). It fails
far from the gate: effect prediction at chance (AUC $0.47$--$0.50$),
$\mathrm{Spearman}\le 0.06$, benefit$-$random $\mathrm{AUC}_{20\text{-}40}=-0.015$, CI
$[-0.035,+0.004]$ pooled. \textbf{E15} goes further and \emph{actively acquires new
information}: before any revision, each agent is asked to name one concrete flaw in its own
answer and rate its confidence --- a genuine meta-cognitive probe, not a re-encoding of the
transcript. This nudges effect-predictability up (AUC $0.47\!\to\!0.55$ pooled, $0.61$ on
GSM8K) --- active information \emph{is} marginally more informative --- yet still falls far
short of the gate, the sign remains wrong ($\mathrm{Spearman}\le 0$), and the probe
controller does not beat a random budget ($\mathrm{AUC}_{20\text{-}40}=-0.018$, CI
$[-0.038,+0.002]$ pooled; significantly negative on BBH). Across three independent,
pre-registered feature families --- cheap geometry, richer semantics, and actively probed
self-diagnosis --- the treatment effect is not identifiable from the pre-intervention state,
despite the large oracle ceiling. Per the pre-registered stop rule we close the method here
and report the gap itself as the contribution; whether a fundamentally richer
intervention-aware signal (a multi-turn probe, or the intervention's own early trace) can
close it is left to future work.

% ---------------------------------------------------------------------------
\begin{figure}[t]
  \centering
  \includegraphics[width=\linewidth]{results/figures/e13_benefit_aware.png}
  \caption{\textbf{E13 --- benefit prediction is the bottleneck, not failure prediction.}
  Policy value (debate accuracy) vs.\ intervention budget for five targeting rules, pooled
  and per benchmark. An \emph{oracle} that targets the true treatment effect (gold, dashed)
  reaches the ceiling by a $\sim\!20\%$ budget ($+16$--$20$ points over never), so there is
  large benefit to capture. But a \emph{learned} CATE controller (blue) --- trained to
  predict which debates the safeguard \emph{helps} --- tracks at or below random: the
  treatment effect is unpredictable from round-0 features ($Y_0$ and $Y_1$ are each
  predictable at AUC $0.73/0.67$, but their difference is not, AUC $0.58$). Failure
  prediction $\neq$ benefit prediction.}
  \label{fig:e13}
\end{figure}
